\section{Results}

\subsection{H-E1: Behavioral Variance Exists}

\begin{table}[h]
\centering
\begin{tabular}{lccc}
\toprule
Metric & Value & Threshold & Status \\
\midrule
Residual Ratio & \textbf{0.6758} & $> 0.05$ & \textbf{PASS} \\
Factor Above Threshold & 13.5$\times$ & --- & --- \\
\bottomrule
\end{tabular}
\caption{Existence hypothesis (H-E1) results.}
\label{tab:he1}
\end{table}

67.6\% of class-wise variance is unexplained by overall accuracy. Vehicle classes (automobile: $\sigma^2 = 0.163$, truck: $\sigma^2 = 0.127$) show highest behavioral variance.

\subsection{H-M1: Weight Statistics Fail}

\begin{table}[h]
\centering
\begin{tabular}{lccc}
\toprule
Metric & Value & Threshold & Status \\
\midrule
Weight $R^2$ & \textbf{-0.08} & --- & --- \\
Baseline $R^2$ & \textbf{0.12} & --- & --- \\
$\Delta R^2$ & \textbf{-0.20} & $> 0$ & \textbf{FAIL} \\
\bottomrule
\end{tabular}
\caption{Mechanism hypothesis (H-M1) results.}
\label{tab:hm1}
\end{table}

Simple weight statistics perform worse than the stratified baseline. Class 4 (deer) shows extreme negative $R^2 = -2.74$, suggesting overfitting on small sample.
